\documentclass[12pt,a4paper]{article}
\usepackage{amsmath, amssymb, amsthm}
\usepackage{graphicx}
\usepackage{booktabs}
\usepackage{hyperref}
\usepackage{xcolor}
\usepackage{algorithm}
\usepackage{algorithmic}
\usepackage{geometry}
\geometry{margin=2.5cm}

\title{\textbf{Multi-Label Zone-Aware Adaptive Heterogeneous GNN\\
for Urban Traffic Flow Prediction under Non-IID Conditions}\\[0.5em]
\large{[Research Plan \& Paper Draft]}}
\author{[Author Name] \\ University of Technology, Ho Chi Minh City}
\date{\today}

\begin{document}
\maketitle
\tableofcontents
\newpage

%% ============================================================
\section{Introduction \& Motivation}
%% ============================================================

Urban traffic prediction is a core problem in Intelligent Transportation Systems (ITS).
Graph Neural Networks (GNNs) have emerged as the dominant approach by modeling road
networks as graphs $\mathcal{G} = (\mathcal{V}, \mathcal{E}, \mathbf{A})$, where nodes
$\mathcal{V}$ represent Traffic Analysis Zones (TAZs), edges $\mathcal{E}$ represent roads, and
$\mathbf{A} \in \mathbb{R}^{N \times N}$ is the adjacency matrix encoding connectivity.

\subsection{The Core Problem: Semantic Blindness of Adjacency Matrices}

\begin{quote}
\textit{``The adjacency matrix encodes \textbf{topology} but is blind to
\textbf{semantics}. Two adjacent nodes may have fundamentally different
traffic patterns if one lies in an industrial zone and the other near a school---
yet the adjacency matrix treats them identically.''}
\end{quote}

Specifically, consider an intersection $v$ located at the boundary of both an
\textbf{industrial zone} and a \textbf{school zone} in Ho Chi Minh City (e.g., near
High Tech Park / VNU-HCM). This node exhibits:

\begin{itemize}
    \item \textbf{Bimodal rush-hour peaks}: both 7:00--8:30 (school) and
          17:00--18:30 (industrial shift end)
    \item \textbf{Weekday/weekend asymmetry}: industrial traffic drops on
          weekends while school traffic drops on school holidays
    \item \textbf{Heterogeneous congestion patterns} not shared by purely
          industrial or purely residential nodes
\end{itemize}

A standard adjacency matrix $\mathbf{A}[v][u]$ only stores whether $v$ and $u$
are connected and with what weight (distance, travel time). It \textbf{cannot}
represent the multi-zone semantic character of node $v$, leading to:

\begin{enumerate}
    \item Incorrect neighborhood aggregation: nodes with different zone compositions
          are treated as equivalent message-passing sources.
    \item Failure to model Non-IID node distributions: the traffic distribution
          $P_v(x)$ differs systematically across zone types.
    \item Loss of interpretability: the model cannot reason about \textit{why}
          a node has anomalous patterns.
\end{enumerate}

%% ============================================================
\section{Related Work}
%% ============================================================

\subsection{GNN-based Traffic Prediction}
DCRNN~\cite{li2018diffusion} introduced diffusion convolution on road graphs.
Graph WaveNet~\cite{wu2019graph} proposed adaptive adjacency learning via
node embeddings. ASTGNN~\cite{guo2021learning} and STGCN~\cite{yu2018spatio}
combined spatial and temporal convolutions. Our baseline AH-GNN extends
Graph WaveNet with time-varying weight matrices.

\textbf{Gap:} None of these works incorporate semantic zone information into
the graph representation or weight generation.

\subsection{Heterogeneous Graph Networks}
HAN~\cite{wang2019heterogeneous} and HeCo~\cite{wang2021self} handle graphs with
multiple node/edge types. SimpleHGN~\cite{lv2021we} scales heterogeneous GNN to
large graphs.

\textbf{Gap:} These methods require \textit{discrete} node types. Urban zones are
\textit{multi-label} (a node can simultaneously belong to commercial, residential,
and school zones), which existing heterogeneous GNNs cannot handle natively.

\subsection{Urban Computing with Zone Context}
Urban2Vec~\cite{wang2020urban2vec} learns neighborhood representations from
POI and street-view data. ZoST~\cite{} incorporates zone semantic for crime prediction.

\textbf{Gap:} No work applies multi-label zone-aware representations to
graph-based traffic flow prediction.

%% ============================================================
\section{Problem Formulation}
%% ============================================================

\subsection{Graph Definition}

We define the urban road network as an attributed multi-source graph:

\begin{equation}
\mathcal{G} = (\mathcal{V},\ \mathcal{E},\ \mathbf{A},\ \mathbf{X}_t,\ \mathbf{Z})
\end{equation}

where:
\begin{itemize}
    \item $\mathcal{V} = \{v_1, \ldots, v_N\}$: set of $N$ traffic nodes
          (Traffic Analysis Zones (TAZs) / districts)
    \item $\mathcal{E}$: road segments connecting nodes
    \item $\mathbf{A} \in \mathbb{R}^{N \times N}$: \textbf{static adjacency matrix}
          derived from OSRM routing data, where
          $A_{ij} = 1/\overline{d}_{ij}$ and $\overline{d}_{ij}$ is the
          mean travel duration from $v_i$ to $v_j$
    \item $\mathbf{X}_t \in \mathbb{R}^{N \times F}$: \textbf{dynamic node features}
          at time $t$, from TomTom API:
          $\mathbf{x}_t^{(v)} = [\text{congestion\_ratio},\ \text{travel\_time},\
          \text{delay},\ \text{free\_flow\_ratio}]$
    \item $\mathbf{Z} \in \{0,1\}^{N \times K}$: \textbf{multi-label zone matrix}
          from OpenStreetMap, where $Z_{vk} = 1$ if node $v$ belongs to zone
          type $k \in \{\text{commercial, residential, industrial, school,}$
          $\text{hospital, university, transport, park}\}$
\end{itemize}

\subsection{Why Adjacency Matrix Is Insufficient: Formal Statement}

\begin{definition}[Semantic Neighborhood]
For node $v$, its \textit{semantic neighborhood} is:
\[
\mathcal{N}_s(v) = \{u \in \mathcal{V} : \mathbf{z}^{(v)} \cdot \mathbf{z}^{(u)} > 0 \}
\]
i.e., nodes sharing at least one zone type.
\end{definition}

\begin{proposition}
The adjacency matrix $\mathbf{A}$ cannot distinguish between nodes with
different zone compositions when $A_{ij} = A_{ik}$ for $j,k$ with
$\mathbf{z}^{(j)} \neq \mathbf{z}^{(k)}$.
\end{proposition}

\begin{proof}
Standard graph convolution aggregates:
$\mathbf{h}_v^{(l+1)} = \sigma\!\left(\sum_{u \in \mathcal{N}(v)} \frac{A_{vu}}{d_v}
\mathbf{h}_u^{(l)} \mathbf{W}\right)$
where $\mathbf{W}$ is shared across all nodes. If $A_{vj} = A_{vk}$, nodes
$j$ and $k$ contribute equally to $v$'s update regardless of their zone labels.
\end{proof}

\subsection{Task Definition}

Given historical observations $\{\mathbf{X}_{t-T+1}, \ldots, \mathbf{X}_t\}$,
static adjacency $\mathbf{A}$, and zone labels $\mathbf{Z}$, predict:

\begin{equation}
\hat{\mathbf{X}}_{t+1} = f_\theta\!\left(\mathbf{X}_{t-T+1:t},\ \mathbf{A},\ \mathbf{Z}\right)
\end{equation}

%% ============================================================
\section{Proposed Method: Zone-Aware AH-GNN}
%% ============================================================

\subsection{Overview}

Our model introduces three key innovations over standard AH-GNN:

\begin{enumerate}
    \item \textbf{Zone Embedding Module}: encodes multi-label zone vector
          $\mathbf{z}^{(v)}$ into a dense representation
    \item \textbf{Zone-Modulated Weight Generation}: node-specific transformation
          weights $\mathbf{W}_v$ conditioned on both spatial embedding and zone embedding
    \item \textbf{Zone-Biased Adaptive Adjacency}: learned adjacency that rewards
          connections between semantically similar nodes
\end{enumerate}

\subsection{Module 1: Zone Embedding}

For each node $v$, we encode its multi-hot zone vector
$\mathbf{z}^{(v)} \in \{0,1\}^K$ into a dense zone embedding:

\begin{equation}
\tilde{\mathbf{z}}^{(v)} = \text{MLP}_{\text{zone}}\!\left(\mathbf{z}^{(v)}\right)
= \sigma\!\left(\mathbf{W}_2 \cdot \text{ReLU}(\mathbf{W}_1 \mathbf{z}^{(v)} + \mathbf{b}_1) + \mathbf{b}_2\right)
\in \mathbb{R}^{d_z}
\end{equation}

This embedding captures semantic similarity between zone compositions:
nodes with $\mathbf{z}^{(v)} = [1,0,1,0,\ldots]$ (commercial + industrial) will
have embeddings closer to nodes with $\mathbf{z}^{(v)} = [1,0,0,0,\ldots]$
(commercial only) than to $\mathbf{z}^{(v)} = [0,0,0,1,\ldots]$ (school only).

\subsection{Module 2: Zone-Aware Adaptive Adjacency}

Following Graph WaveNet, we define a learnable node embedding
$\mathbf{E} \in \mathbb{R}^{N \times d_e}$. We augment it with zone information:

\begin{equation}
\tilde{\mathbf{E}}_v = \mathbf{E}_v + \mathbf{W}_{\text{proj}} \tilde{\mathbf{z}}^{(v)}
\end{equation}

The time-varying adaptive adjacency is then:
\begin{equation}
\tilde{\mathbf{A}}_t = \text{Softmax}\!\left(\text{ReLU}\!\left(
\tilde{\mathbf{E}} \cdot \mathbf{W}_t \cdot \tilde{\mathbf{E}}^\top
\right)\right)
\end{equation}

where $\mathbf{W}_t \in \mathbb{R}^{d_e \times d_e}$ is a time-label-specific
weight matrix (one per time label: night, morning rush, evening rush, normal).

\textbf{Combined adjacency} with static OSRM data:
\begin{equation}
\mathbf{A}_{\text{final}} = \alpha \cdot \tilde{\mathbf{A}}_t + (1-\alpha) \cdot \mathbf{A}_{\text{OSRM}}
\end{equation}

where $\alpha \in [0,1]$ is a learnable scalar. This allows the model to
\textit{blend} learned semantic adjacency with observed routing structure.

\subsection{Module 3: Zone-Modulated Node-Specific Weights}

Standard GNN uses a single shared weight $\mathbf{W}$ for all nodes.
We instead generate \textit{node-specific} weights conditioned on zone:

\begin{equation}
\mathbf{W}_v = \text{reshape}\!\left(
\mathbf{W}_{\text{gen}} \cdot \text{concat}(\mathbf{E}_v,\ \tilde{\mathbf{z}}^{(v)})
\right) \in \mathbb{R}^{d_{\text{in}} \times d_{\text{out}}}
\end{equation}

\begin{equation}
\mathbf{b}_v = \mathbf{b}_{\text{gen}} \cdot \text{concat}(\mathbf{E}_v,\ \tilde{\mathbf{z}}^{(v)})
\in \mathbb{R}^{d_{\text{out}}}
\end{equation}

This means \textit{nodes with different zone combinations use different
transformation weights}, solving the proposition stated in Section 3.2.

\subsection{Forward Pass}

One layer of Zone-Aware Adaptive Graph Convolution:

\begin{equation}
\mathbf{H}^{(l+1)}_v = \text{ReLU}\!\left(
\sum_{u \in \mathcal{V}} [\mathbf{A}_{\text{final}}]_{vu} \cdot
\mathbf{H}^{(l)}_u \cdot \mathbf{W}_v + \mathbf{b}_v
\right)
\end{equation}

The input to the first layer fuses temporal features with zone embedding:
\begin{equation}
\mathbf{H}^{(0)}_v = \text{concat}\!\left(\mathbf{x}_t^{(v)},\ \tilde{\mathbf{z}}^{(v)}\right)
\end{equation}

\subsection{Full Model Architecture}

\begin{equation}
\hat{\mathbf{X}}_{t+1} = \text{FC}\!\left(
\text{ZoneAwareAHGNN}\!\left(\mathbf{X}_t,\ \mathbf{Z},\ \tau_t,\ \mathbf{A}_{\text{OSRM}}\right)
\right)
\end{equation}

\noindent\textbf{Parameter count comparison:}

\begin{table}[h]
\centering
\begin{tabular}{lcc}
\toprule
\textbf{Component} & \textbf{AH-GNN (baseline)} & \textbf{Proposed} \\
\midrule
Node embedding $\mathbf{E}$ & $N \times d_e$ & $N \times d_e$ \\
Zone embedding MLP & --- & $K \times 2d_z + 2d_z \times d_z$ \\
Time weights $\mathbf{W}_t$ & $|\mathcal{T}| \times d_e^2$ & $|\mathcal{T}| \times d_e^2$ \\
Weight generator & $d_e \times (d_{\text{in}} \cdot d_{\text{out}})$ & $(d_e+d_z) \times (d_{\text{in}} \cdot d_{\text{out}})$ \\
Zone proj $\mathbf{W}_{\text{proj}}$ & --- & $d_z \times d_e$ \\
\bottomrule
\end{tabular}
\caption{Parameter comparison. Additional zone-related parameters are minimal.}
\end{table}

%% ============================================================
\section{Data Pipeline: Combining 3 Sources}
%% ============================================================

\subsection{Source 1 — Static Adjacency Matrix (OSRM)}

\textbf{Source:} \texttt{hcm\_osrm\_dataset.csv} (20,379 rows, 17 nodes)\\
\textbf{API:} \texttt{router.project-osrm.org} (free, based on OpenStreetMap road network)\\
\textbf{Usage:} Build $\mathbf{A}_{\text{OSRM}}$

\begin{equation}
A_{ij}^{\text{OSRM}} = \frac{1}{\overline{d}_{ij}}, \quad
\overline{d}_{ij} = \frac{1}{|S_{ij}|} \sum_{s \in S_{ij}} d_{ij}^{(s)}
\end{equation}

where $S_{ij}$ is the set of sampled routes from node $i$ to $j$.

\textbf{Limitation:} \texttt{duration\_s} in OSRM is \textit{static routing time},
not real-time traffic. It does not reflect congestion.

\subsection{Source 2 — Dynamic Traffic Features (TomTom)}

\textbf{Source:} TomTom Routing API v1 (\texttt{tomtom\_collector.py})\\
\textbf{API:} \texttt{api.tomtom.com/routing/1/calculateRoute} with \texttt{traffic=true}\\
\textbf{Coverage:} 10 nodes in central Ho Chi Minh City, 90 directed pairs\\
\textbf{Usage:} Build $\mathbf{X}_t$ (node feature matrix per snapshot)

Per edge $(i \to j)$ at time $t$, TomTom returns:
\begin{itemize}
    \item \texttt{travel\_time\_s} $= tt_{ij}^t$: actual travel time with traffic
    \item \texttt{free\_flow\_time\_s} $= ff_{ij}^t$: ideal time without traffic
    \item \texttt{traffic\_delay\_s} $= dl_{ij}^t$: delay due to congestion
    \item \texttt{congestion\_ratio} $= tt_{ij}^t / ff_{ij}^t \geq 1$
\end{itemize}

\textbf{Edge-to-Node aggregation} (pivot edge features to node features):
\begin{equation}
\mathbf{x}_t^{(v)} = \frac{1}{|\mathcal{N}^+(v)|}
\sum_{u \in \mathcal{N}^+(v)}
\left[\text{congestion}_{vu}^t,\ \text{delay}_{vu}^t,\ \text{tt}_{vu}^t,\ \text{ff\_ratio}_{vu}^t \right]
\end{equation}

where $\mathcal{N}^+(v)$ is the set of outgoing neighbors of $v$.

\textbf{Time labels} for $\mathbf{W}_t$ selection:
\begin{equation}
\tau_t = \begin{cases}
0 & \text{night } (0\text{h}–6\text{h}) \\
1 & \text{morning rush } (7\text{h}–9\text{h}) \\
2 & \text{evening rush } (16\text{h}–19\text{h}) \\
3 & \text{normal } (\text{otherwise})
\end{cases}
\end{equation}

\subsection{Source 3 — Zone Labels (OpenStreetMap / Overpass API)}

\textbf{Source:} Overpass API (\texttt{overpass-api.de/api/interpreter})\\
\textbf{Usage:} Build $\mathbf{Z} \in \{0,1\}^{N \times K}$ (static, collected once)

For each node $v$ at coordinates $(lat_v, lon_v)$, query all OSM features within
radius $r = 800$m:

\begin{align}
Z_{v,\text{commercial}} &= \mathbf{1}[\exists \text{ way with landuse=commercial in } B(v, r)] \\
Z_{v,\text{industrial}}  &= \mathbf{1}[\exists \text{ way with landuse=industrial in } B(v, r)] \\
Z_{v,\text{school}}      &= \mathbf{1}[\exists \text{ way with amenity=school in } B(v, r)] \\
Z_{v,\text{university}}  &= \mathbf{1}[\exists \text{ way with amenity=university in } B(v, r)] \\
Z_{v,\text{hospital}}    &= \mathbf{1}[\exists \text{ way with amenity=hospital in } B(v, r)] \\
Z_{v,\text{transport}}   &= \mathbf{1}[\exists \text{ aeroway or landuse=transportation in } B(v, r)] \\
Z_{v,\text{residential}} &= \mathbf{1}[\exists \text{ way with landuse=residential in } B(v, r)] \\
Z_{v,\text{park}}        &= \mathbf{1}[\exists \text{ way with leisure=park in } B(v, r)]
\end{align}

\textbf{Expected zone labels for HCM nodes:}

\begin{table}[h]
\centering
\small
\begin{tabular}{lccccccc}
\toprule
\textbf{Node} & \textbf{comm.} & \textbf{resid.} & \textbf{indust.} & \textbf{school} & \textbf{univ.} & \textbf{hosp.} & \textbf{transp.} \\
\midrule
Ben Thanh Market    & 1 & 0 & 0 & 0 & 0 & 0 & 0 \\
District 1          & 1 & 1 & 0 & 0 & 0 & 1 & 0 \\
Tan Son Nhat Airport& 1 & 0 & 0 & 0 & 0 & 0 & 1 \\
High Tech Park      & 0 & 0 & 1 & 0 & 0 & 0 & 0 \\
VNU HCM             & 0 & 1 & 0 & 1 & 1 & 0 & 0 \\
Suoi Tien           & 1 & 0 & 0 & 0 & 0 & 0 & 0 \\
\bottomrule
\end{tabular}
\caption{Expected multi-label zone matrix for selected HCM nodes.
         VNU HCM has \textbf{two active zone labels} (school + university),
         which is the case the adjacency matrix cannot represent.}
\end{table}

%% ============================================================
\section{Experiments}
%% ============================================================

\subsection{Dataset Summary}

\begin{table}[h]
\centering
\begin{tabular}{llll}
\toprule
\textbf{Source} & \textbf{Role in Model} & \textbf{Current Status} & \textbf{Size} \\
\midrule
OSRM & $\mathbf{A}_{\text{OSRM}}$ (static adj.) & \textcolor{green}{Available} & 20K rows, 17 nodes \\
TomTom & $\mathbf{X}_t$ (dynamic features) & \textcolor{orange}{Need to collect} & Target: 2 weeks \\
OSM/OVP & $\mathbf{Z}$ (zone labels) & \textcolor{orange}{Need to query} & 1 query per node \\
\bottomrule
\end{tabular}
\caption{Dataset status as of June 2026.}
\end{table}

\subsection{Baselines}

\begin{enumerate}
    \item \textbf{DCRNN}: Diffusion Convolutional RNN (Li et al., 2018)
    \item \textbf{Graph WaveNet}: Adaptive adjacency without zone (Wu et al., 2019)
    \item \textbf{AH-GNN} (current): Time-varying adaptive adjacency, no zone labels
    \item \textbf{AH-GNN + Zone (concat)}: Zone features naively concatenated
    \item \textbf{Proposed (Zone-Aware AH-GNN)}: Full model with zone embedding
          + zone-modulated weights + zone-biased adjacency
\end{enumerate}

\subsection{Ablation Study}

\begin{table}[h]
\centering
\begin{tabular}{lccc}
\toprule
\textbf{Variant} & \textbf{Zone Embed} & \textbf{Zone-Weight} & \textbf{Zone-Adj} \\
\midrule
AH-GNN (baseline) & \xmark & \xmark & \xmark \\
+ Zone Concat      & \cmark & \xmark & \xmark \\
+ Zone Weight      & \cmark & \cmark & \xmark \\
Full (Proposed)    & \cmark & \cmark & \cmark \\
\bottomrule
\end{tabular}
\caption{Ablation variants to isolate contribution of each zone-aware component.}
\end{table}

\subsection{Evaluation Metrics}

\begin{itemize}
    \item MAE: $\frac{1}{N}\sum_v |\hat{x}_v - x_v|$
    \item RMSE: $\sqrt{\frac{1}{N}\sum_v (\hat{x}_v - x_v)^2}$
    \item MAPE: $\frac{1}{N}\sum_v \frac{|\hat{x}_v - x_v|}{x_v} \times 100\%$
    \item \textbf{Zone-Stratified MAE}: compute MAE separately for each zone type
          to show improvement specifically for multi-zone nodes
\end{itemize}

\subsection{Key Experiment: Distribution Heterogeneity}

To motivate the problem (motivation section), compute Jensen-Shannon Divergence
between node pairs with different zone combinations:

\begin{equation}
\text{JSD}(P_v \| P_u) = \frac{1}{2} D_{\text{KL}}(P_v \| M) + \frac{1}{2} D_{\text{KL}}(P_u \| M),
\quad M = \frac{P_v + P_u}{2}
\end{equation}

\textbf{Hypothesis:} Nodes with disjoint zone labels (e.g., industrial vs.\ school)
will have higher JSD than nodes with overlapping zones.

%% ============================================================
\section{Implementation Plan}
%% ============================================================

\begin{enumerate}
    \item \textbf{Week 1}: Run Overpass API queries for all 17 nodes → \texttt{node\_zone\_labels.csv}
    \item \textbf{Week 1--2}: Collect TomTom data with \texttt{tomtom\_collector.py}
          for at least 7 days (capturing full weekly pattern including rush hours)
    \item \textbf{Week 2}: Build adjacency matrix from OSRM, pivot TomTom edges to
          node features, compute JSD motivation statistics
    \item \textbf{Week 3}: Implement and train Zone-Aware AH-GNN, run ablations
    \item \textbf{Week 4}: Write paper (Introduction through Experiments)
\end{enumerate}

%% ============================================================
\section{Expected Contributions}
%% ============================================================

\begin{enumerate}
    \item \textbf{Problem formulation}: first formal definition of multi-label
          zone-aware traffic graph with proof of adjacency matrix insufficiency
    \item \textbf{Architecture}: Zone Embedding + Zone-Modulated Weight Generation
          + Zone-Biased Adaptive Adjacency
    \item \textbf{Dataset}: HCM-Zone — first public dataset combining OSRM routing,
          TomTom real-time traffic, and OSM multi-label zone annotations for
          Ho Chi Minh City
    \item \textbf{Empirical}: zone-stratified evaluation demonstrating particular
          gains at multi-zone nodes
\end{enumerate}

%% ============================================================
\begin{thebibliography}{99}
\bibitem{li2018diffusion} Li, Y., et al. (2018). Diffusion Convolutional Recurrent Neural
Network. \textit{ICLR}.
\bibitem{wu2019graph} Wu, Z., et al. (2019). Graph WaveNet. \textit{IJCAI}.
\bibitem{yu2018spatio} Yu, B., et al. (2018). Spatio-Temporal Graph Convolutional Networks.
\textit{IJCAI}.
\bibitem{guo2021learning} Guo, S., et al. (2021). Learning Dynamics and Heterogeneity of
Spatial-Temporal Graph Data for Traffic Forecasting. \textit{IEEE TKDE}.
\bibitem{wang2019heterogeneous} Wang, X., et al. (2019). Heterogeneous Graph Attention
Network. \textit{WWW}.
\bibitem{wang2021self} Wang, X., et al. (2021). Self-supervised Heterogeneous Graph
Pre-training. \textit{NeurIPS}.
\bibitem{lv2021we} Lv, Q., et al. (2021). Are We Really Making Much Progress? \textit{KDD}.
\bibitem{wang2020urban2vec} Wang, Z., et al. (2020). Urban2Vec. \textit{AAAI}.
\end{thebibliography}

\end{document}
